\documentclass[11pt,letterpaper]{article}

% ==========================================================================
% SPSE4 v0.1 -- System Description and Design Rationale
% Andrew Dougherty, FRDCSA / FRKCSA Project, April 2026
%
% This LaTeX source is hand-authored to match the master Org-mode document
% (spse4.org) exactly.  If you edit the Org file, regenerate the PDF either
% by running the supplied build.sh (which uses this .tex as the master for
% typesetting) or by hand-updating this file.
% ==========================================================================

\usepackage[margin=1in]{geometry}
\usepackage[activate={true,nocompatibility},protrusion=true,expansion=false]{microtype}
\usepackage{booktabs}
\usepackage{enumitem}
\usepackage{xcolor}
\usepackage{titlesec}
\usepackage[most]{tcolorbox}
\usepackage{natbib}
\usepackage{hyperref}
\usepackage{doi}
\usepackage{listings}
\usepackage{fancyhdr}

\hypersetup{
  colorlinks=true,
  linkcolor=blue!50!black,
  urlcolor=blue!50!black,
  citecolor=blue!50!black
}

\titleformat{\section}{\normalfont\Large\bfseries}{\thesection}{1em}{}
\titleformat{\subsection}{\normalfont\large\bfseries}{\thesubsection}{1em}{}
\setlist{nosep,leftmargin=*,topsep=2pt}

\definecolor{codebg}{RGB}{248,248,248}
\definecolor{coderule}{RGB}{220,220,220}

\lstset{%
  basicstyle=\ttfamily\small,
  breaklines=true,
  backgroundcolor=\color{codebg},
  frame=single,
  rulecolor=\color{coderule},
  framesep=4pt,
  xleftmargin=6pt,
  xrightmargin=6pt,
  aboveskip=6pt,
  belowskip=6pt,
  columns=fullflexible,
  keepspaces=true,
  showstringspaces=false
}

\pagestyle{fancy}
\fancyhf{}
\lhead{\small SPSE4 v0.1}
\rhead{\small FRDCSA / FRKCSA}
\cfoot{\thepage}
\renewcommand{\headrulewidth}{0.2pt}

\setlength{\parskip}{4pt plus 1pt minus 1pt}
\setlength{\parindent}{0pt}

\title{\textbf{SPSE4: A Multi-Relation Graph-Based Planning-Domain Editor for SWI-Prolog}\\
       \large System Description and Design Rationale for v0.1}
\author{Andrew Dougherty\\
        \small FRDCSA / FRKCSA Project\\
        \small \texttt{adougher9@yahoo.com}}
\date{May 2, 2026}

\begin{document}
\maketitle
\thispagestyle{fancy}

\begin{abstract}
\noindent SPSE4 is a multi-relation graph-based planning-domain editor for
SWI-Prolog, implemented as five composable Prolog packs totaling
approximately 3,900 lines with 79 unit tests. It succeeds SPSE2 (2011,
Perl/Tk) and SPSE2-Formalog (SWI-Prolog bridge) by consolidating the task
ontology, temporal reasoning, planning-language I/O, scheduling, and
reactive execution into a single process with a single data model, rooted in
microtheories rather than Perl UniLang round-trips. The system integrates
Allen's interval algebra with CLP(FD) for schedule feasibility, persistent
microtheories with per-context access control and audit, a full PDDL/HDDL
parser and emitter with a planner-capability match function, and a
real-datetime scheduler with critical-path analysis and broadcast-based
reactive triggers. This report describes the v0.1 architecture, the
pack-by-pack decomposition, the key design decisions with their rationale,
a worked end-to-end example driving the FRDCSA autopackager's priority
queue, and the roadmap to v0.2.
\end{abstract}

\begin{center}
\fbox{\parbox{0.92\linewidth}{\centering\small
\textbf{Draft.}\ This document is a working draft under active revision.
Technical descriptions reflect the v0.1 codebase as of the date above, but
prose, citations, and architectural framing may contain inaccuracies and
are subject to change. Please consult the GitHub repository at
\url{https://github.com/aindilis/spse4} for the current source and the
latest revision of this document. Feedback is welcome.
}}
\end{center}
\vspace{0.5em}

\tableofcontents
\vspace{1em}

% ==========================================================================
\section{Introduction}
% ==========================================================================

The Shared Priority System Editor v2 (SPSE2) was described in the 2011
FRDCSA technical report \emph{Temporal Planning and Inferencing for
Personal Task Management with SPSE2} \citep{spse2-2011}. Its
architecture tied together a Perl/Tk GUI, the FRDCSA UniLang messaging
layer, the FreeKBS2 knowledge-base server, the Verber PDDL translator,
and the LPG-TD planner. A companion module, SPSE2-Formalog, bridged
portions of that world to SWI-Prolog so that knowledge-engineering and
inference could happen where the Prolog infrastructure already was.

Over fifteen years of use, two things became clear. First, the round-trip
architecture---Perl emitting UniLang messages to a FreeKBS2 image, Prolog
pulling facts by \emph{kquery} and flattening them into \texttt{holds/2},
Verber translating back out to PDDL, and a planner returning a plan that got
re-read into the GUI---was a source of both latency and fragility. Second,
the core data model (typed goals, typed binary edges, per-context
assertion) was sound and worth preserving. What needed rebuilding was the
infrastructure around it, not the ontology.

SPSE4 is that rebuild. It is not a port of SPSE2 to Prolog; it is a
first-principles redesign of a planning-domain editor as a set of
composable Prolog packs, any of which can be used standalone. The
target audience is threefold: (1) the author's own daily planning and
scheduling needs via the Free Life Planner (FLP); (2) the FRDCSA
\citep{frdcsa-project} subsystem ecosystem, including the
autopackager's work queue; and (3) the wider SWI-Prolog and
AI-planning communities, for whom \texttt{pack-allen},
\texttt{pack-mt-store}, and \texttt{pack-pddl} solve problems with no
other clean SWI-Prolog solutions.

This v0.1 release ships the five foundational packs plus a live
demonstration seeded with the autopackager's actual priority queue. It does
not ship the collaboration server, the browser client, or the Emacs mode;
those belong to v0.2, which is sketched in Section~\ref{sec:roadmap}.

% ==========================================================================
\section{Background: What SPSE2 Was, and What Needed Rebuilding}
% ==========================================================================

The SPSE2 technical report was prepared as a submission to the ICAPS
2011 Workshop on Scheduling and Planning Applications but was not
accepted for publication; it has remained available as a project
document on the FRDCSA website. It describes a system with three
conceptually separate layers entangled through the FRDCSA middleware:

\begin{enumerate}
\item An \emph{ontology} of goals as graph nodes, with unary predicates
  marking status (\texttt{complete}, \texttt{cancelled}, \texttt{deleted},
  \texttt{showstopper}, \texttt{ridiculous}, \texttt{obsoleted},
  \texttt{rejected}, \texttt{skipped}) and binary predicates marking
  structural relations (\texttt{depends}, \texttt{provides}, \texttt{eases},
  \texttt{prefer}). Per-goal metadata included \texttt{has-NL}
  (natural-language label), \texttt{has-source}, \texttt{asserter},
  \texttt{start-date}/\texttt{end-date}, \texttt{costs}, \texttt{earns},
  \texttt{severity}, \texttt{event-duration}, \texttt{has-formalization},
  \texttt{has-parent}, \texttt{involves}, \texttt{ethicality-concern}, and
  \texttt{hard-time-constraints}.

\item A \emph{PDDL bridge}: SPSE2 $\to$ Verber $\to$ PDDL $\to$ LPG-TD $\to$
  plan, with calendar events (ICS, Google Calendar) translated into PDDL
  2.2 timed initial literals and a durative \texttt{Complete} action gated
  by an \texttt{over all ((not has-time-constraints ?e) or possible(?e))}
  condition.

\item A \emph{context system}: everything lived inside
  \texttt{holds(Context, Fact)}, with the contexts serving as microtheories
  in the Cyc sense---per-domain scopes the user could switch between.
\end{enumerate}

The SPSE2-Formalog layer read contexts via shelling out to the Perl
\texttt{spse2} binary, pulled facts via \texttt{kquery} into
\texttt{holds/2} assertions, and exposed accessors like \texttt{goal/1},
\texttt{depends/2}, \texttt{has-NL/2} that flattened the context prefix \citep{guha-lenat-cyc-1990}.
The actual planning, GUI, and calendar logic remained in Perl/Tk and Emacs
Lisp.

Three properties of that architecture motivated the rebuild:

\begin{itemize}
\item \emph{No persistence beyond the Perl image.} The goal graph lived in
  a running Perl process's memory, re-asserted from dumps on startup.
  Concurrent access was impossible; the image was the bottleneck.
\item \emph{Scheduling round-trips had to go through PDDL.} Even simple
  ``when can this happen given these constraints'' queries involved PDDL
  generation and an external planner. This made sense in 2011 when
  Prolog's CLP(FD) was less mature; in 2026 it is unnecessary overhead for
  common cases.
\item \emph{The context system was ad-hoc.} \texttt{holds/2} worked, but
  there was no ACL, no audit trail, and no genuine inheritance in the Cyc
  \texttt{genlMt} sense.
\end{itemize}

SPSE4 addresses all three by (1) persisting microtheories through a
dedicated store pack, (2) providing native CLP(FD)-based scheduling as the
common case while keeping PDDL emission available for the hard cases, and
(3) making microtheories first-class objects with owners, visibility,
inheritance, and audit.

% ==========================================================================
\section{Design Rationale: Six Decisions That Shaped v0.1}
\label{sec:rationale}
% ==========================================================================

Before the architecture is described, six design decisions taken during the
v0.1 design session are worth stating explicitly, because they determined
what v0.1 \emph{is} and what it deliberately \emph{is not}.

\subsection{Native Prolog everywhere, no Perl round-trips}

The entire SPSE4 runtime is SWI-Prolog. The SPSE2 legacy dump format is
supported as an \emph{importer}, not as a live bridge. A Perl installation
is not required to use SPSE4. This decision follows from the consolidation
argument in \S2: the main pain in SPSE2 was the seams between components,
and the simplest way to eliminate seams is to eliminate components.

\subsection{Microtheories as first-class, persistent, ACL-aware objects}

The \texttt{holds(Context, Fact)} pattern from SPSE2 is replaced by
\texttt{ist(Mt, Fact)} in \texttt{pack-mt-store} (naming follows Guha's
CycL convention \citep{guha-lenat-cyc-1990}), with \texttt{genlMt(Sub, Super)} for inheritance,
per-microtheory owner and visibility, and an append-only audit log of every
assertion and retraction. This is strictly more structure than SPSE2
provided, but the extra machinery is what enables the collaboration server
planned for v0.2: you cannot broadcast diffs safely without knowing who
asserted what, when, under whose authority.

\subsection{Scheduling is CLP(FD)-native; PDDL is an exporter}

Allen's thirteen interval relations \citep{allen-1983} are implemented directly over CLP(FD)
\citep{triska-clpfd-2012} finite-domain variables in \texttt{pack-allen}. Common scheduling
questions---``is this set of relations consistent?'', ``find concrete
start/end times for these intervals''---are answered by the constraint
solver without any PDDL round trip. \texttt{pack-pddl} emits (and parses)
PDDL when the user explicitly wants to hand a domain and problem to an
external planner, but this is never on the critical path for daily use.

\subsection{Native Prolog behavior-tree executor, not BT.CPP}

The reactive execution layer in v0.2 will be a native Prolog behavior-tree
interpreter, not a binding to BT.CPP. Rationale: the FLP use case is
cognitive-prosthetic task execution, not robotics; BT.CPP's runtime
properties (real-time guarantees, C++ callable leaves) are not needed, and
embedding Prolog into a BT.CPP tick loop is architecturally awkward. If a
robotics use case ever emerges, BT.CPP XML is a one-way exporter
target---no runtime dependency on the other direction.

\subsection{CycL is emitted as text, not called live}

There is no live SWI-Prolog $\leftrightarrow$ Cyc client in v0.1.
\texttt{pack-spse4-core} emits \texttt{.cycl} files using the vocabulary
from the SPSE2 Logical Operators work (\texttt{\#\$isa},
\texttt{\#\$SPSE2Mt}, \texttt{(am '(\ldots) Mt)} syntax). When a working
Prolog-to-Cyc bridge exists (Doug Miles's PrologCyc or another
collaborator's reimplementation), it can consume these files. The system
builds against a stable text format and does not depend on a live Cyc
image.

\subsection{Each pack is standalone-useful}

The five packs are designed so that each one is worth installing by itself.
\texttt{pack-allen} is a complete Allen interval algebra implementation
for anyone doing temporal reasoning in SWI-Prolog; \texttt{pack-mt-store}
is a persistent microtheory store for anyone doing KR; \texttt{pack-pddl}
is a PDDL parser and emitter for anyone working with AI planners. The
SPSE4-specific packs sit on top of these, but the foundation packs do not
know or care about SPSE4. This increases the blast radius of bug reports
(anyone using \texttt{pack-pddl} may find bugs in it, which benefits
SPSE4) and opens a straightforward contribution path for the wider
community.

% ==========================================================================
\section{Architecture Overview}
\label{sec:arch}
% ==========================================================================

The v0.1 architecture is a three-layer dependency graph:

\begin{center}
\begin{tabular}{@{}l@{}}
\texttt{pack-spse4-scheduler} \quad (datetimes, critical path, reactive triggers)\\
\quad $\downarrow$ depends on $\downarrow$\\
\texttt{pack-spse4-core} \quad (task ontology, edges, projections, SPSE2 import)\\
\quad $\downarrow$ depends on $\downarrow$\\
\texttt{pack-mt-store} \quad \texttt{pack-allen} \quad \texttt{pack-pddl}\\
\end{tabular}
\end{center}

\begin{itemize}
\item The \emph{foundation layer} (\texttt{pack-allen},
  \texttt{pack-mt-store}, \texttt{pack-pddl}) has no SPSE4-specific
  dependencies. Each pack is independently installable and useful to the
  wider community.
\item The \emph{core layer} (\texttt{pack-spse4-core}) defines the task
  ontology and depends on \texttt{pack-mt-store} (for microtheory storage)
  and \texttt{pack-allen} (for temporal relations on task intervals). It
  also contains the SPSE2 legacy importer and the projection queries that
  filter the graph by relation kind, status, or microtheory membership.
\item The \emph{scheduler layer} (\texttt{pack-spse4-scheduler}) adds
  real-datetime scheduling, critical-path analysis, and broadcast-based
  reactive triggers, depending on \texttt{pack-spse4-core} and
  \texttt{pack-allen}.
\end{itemize}

All packs follow the Prolog coding conventions of \citet{covington-style-2012}
and \citet{bratko-2011}: \texttt{:- module/2} on every file, PlDoc docstrings
on every exported predicate, PlUnit test files colocated in \texttt{t/},
strings for all user-visible text, and atoms for identifiers and enumerated
values. These conventions are also consistent with the broader SWI-Prolog
ecosystem \citep{wielemaker-swipl-2012}.

% ==========================================================================
\section{The Five Packs}
\label{sec:packs}
% ==========================================================================

\subsection{\texttt{pack-allen}: Allen Interval Algebra + CLP(FD) Scheduler}
\label{sec:allen}

Approximately 200 lines, 9 tests. Implements the thirteen primitive binary
relations between closed intervals on a linear timeline \citep{allen-1983}: \texttt{before},
\texttt{meets}, \texttt{overlaps}, \texttt{starts}, \texttt{during},
\texttt{finishes}, \texttt{equals}, and their inverses \texttt{after},
\texttt{met\_by}, \texttt{overlapped\_by}, \texttt{started\_by},
\texttt{contains}, \texttt{finished\_by}.

The module exports:

\begin{itemize}
\item \texttt{allen\_relation/1}, \texttt{allen\_inverse/2},
  \texttt{allen\_compose/3}---the qualitative layer: relation enumeration,
  inverse lookup, and the composition table (what relation holds between
  A and C, given A/B and B/C).
\item \texttt{interval/3}, \texttt{constrain\_interval/3}---interval
  declaration with CLP(FD) start and end variables, and constraint posting
  by qualitative relation.
\item \texttt{schedule/3}---given a list of intervals and a list of
  qualitative constraints, find concrete integer start/end times
  satisfying all constraints.
\item \texttt{relation\_implies/2}---strict-to-loose reasoning (e.g.\
  \texttt{starts} implies ``not after'').
\end{itemize}

A typical call:

\begin{lstlisting}
?- schedule([a, b, c],
            [ constraint(a, before, b),
              constraint(b, overlaps, c) ],
            Solution).
\end{lstlisting}

returns concrete integer timelines for \texttt{a}, \texttt{b}, \texttt{c}
that respect both constraints, or fails if the constraint set is
inconsistent. The \texttt{schedule/3} predicate is the bridge between
qualitative Allen reasoning and concrete integer timelines;
\texttt{pack-spse4-scheduler} layers real datetime arithmetic on top by
treating the integers as Unix-epoch seconds.

\textbf{Why this matters standalone.} There is no other clean
Allen-algebra implementation for SWI-Prolog. The composition table alone
is a reusable contribution; wiring it to CLP(FD) makes it directly useful
for scheduling.

\subsection{\texttt{pack-mt-store}: Persistent Microtheory Store}
\label{sec:mt}

Approximately 300 lines, 10 tests. Implements a microtheory store with
\texttt{ist(Mt, Fact)}, \texttt{genlMt(Sub, Super)},
\texttt{specialization/2}, per-microtheory owner and visibility metadata,
and an append-only audit log. The name and vocabulary follow the CycL
conventions of \citet{guha-lenat-cyc-1990} so the code reads naturally to
anyone familiar with Cyc.

Core predicates: \texttt{mt\_create/2} (create a microtheory with owner
and visibility options); \texttt{mt\_specialize/2} (declare a
\texttt{genlMt} relation, i.e.\ inheritance); \texttt{mt\_assert/2},
\texttt{mt\_retract/2} (assertion and retraction, audit-logged);
\texttt{ist/2} (query facts local to a microtheory);
\texttt{ist\_inherited/2} (query facts visible through the
\texttt{genlMt} lattice); and \texttt{mt\_audit/2} (stream the audit log
for a microtheory).

Example:

\begin{lstlisting}
?- mt_create(general_kb, [owner=andrew, visibility=public]).
?- mt_create(medical_kb, [owner=andrew]).
?- mt_specialize(medical_kb, general_kb).
?- mt_assert(general_kb, organism(andrew, human)).
?- mt_assert(medical_kb, takes(andrew, vitaminD)).
?- ist_inherited(medical_kb, Fact).
Fact = takes(andrew, vitaminD) ;
Fact = organism(andrew, human).
\end{lstlisting}

Vocabulary mapping:

\begin{center}
\begin{tabular}{@{}ll@{}}
\toprule
SPSE2 & SPSE4 \texttt{pack-mt-store} \\
\midrule
\texttt{holds(Context, Fact)} & \texttt{ist(Context, Fact)} \\
context inheritance (ad-hoc)  & \texttt{specialization/2}, \texttt{genlMt/2} \\
\bottomrule
\end{tabular}
\end{center}

\begin{center}
\begin{tabular}{@{}ll@{}}
\toprule
Cyc & SPSE4 \texttt{pack-mt-store} \\
\midrule
\texttt{(ist Mt Fact)}        & \texttt{ist(Mt, Fact)} \\
\texttt{(genlMt Sub Super)}   & \texttt{genlMt(Sub, Super)} \\
\bottomrule
\end{tabular}
\end{center}

v0.1 stores microtheories in-memory with persistence hooks in place but
not yet wired; v0.2 will add the MySQL backend via the planned
\texttt{prolog-mysql-store} integration for true multi-process persistence.

\subsection{\texttt{pack-spse4-core}: Task Ontology and Graph Operations}
\label{sec:core}

Approximately 450 lines, 13 tests. Defines the core data model of SPSE4:
a \emph{priority system} is a labelled directed multigraph stored in one
or more microtheories, with tasks as nodes and typed binary relations as
edges.

All user-visible text (NL labels, descriptions, notes) is stored as
\textbf{strings}; internal identifiers, predicate names, relation kinds,
and enumerated values are \textbf{atoms}. This follows Covington's coding
conventions.

Task properties:

\begin{center}
\small
\begin{tabular}{@{}lp{0.62\linewidth}@{}}
\toprule
Key & Value type \\
\midrule
\texttt{has\_nl}         & string --- natural-language label \\
\texttt{description}     & string --- longer text \\
\texttt{task\_kind}      & atom --- \texttt{primitive} $\vert$ \texttt{compound} $\vert$ \texttt{recurring} $\vert$ \texttt{ongoing} $\vert$ \texttt{milestone} \\
\texttt{overlap\_class}  & atom --- \texttt{exclusive} $\vert$ \texttt{foregroundable} $\vert$ \texttt{backgroundable} $\vert$ \texttt{ambient} $\vert$ \texttt{batched} \\
\texttt{status}          & atom --- \texttt{open} $\vert$ \texttt{in\_progress} $\vert$ \texttt{completed} $\vert$ \texttt{cancelled} $\vert$ \texttt{deleted} $\vert$ \texttt{skipped} $\vert$ \texttt{obsoleted} $\vert$ \texttt{rejected} $\vert$ \texttt{showstopper} $\vert$ \texttt{ridiculous} $\vert$ \texttt{habitual} \\
\texttt{duration}        & term --- \texttt{fixed(N)} $\vert$ \texttt{range(Lo,Hi)} $\vert$ \texttt{distribution(K,Ps)} $\vert$ \texttt{effort(Total)} \\
\texttt{earliest\_start} & datetime term \\
\texttt{latest\_finish}  & datetime term \\
\texttt{needs\_resource} & atom --- may repeat \\
\texttt{costs}, \texttt{earns} & term --- e.g.\ \texttt{dollars(30)} \\
\texttt{recurrence}      & string --- RFC 5545 RRULE \\
\texttt{triggers\_when}  & callable goal \\
\texttt{source}          & arbitrary term (e.g.\ \texttt{spse2(Ctx)} to track provenance) \\
\bottomrule
\end{tabular}
\end{center}

Edges carry an explicit \texttt{relation\_kind} discriminator so the same
pair of tasks can stand in multiple relations simultaneously:

\begin{itemize}
\item \texttt{depends}---A must complete before B can start (causal prerequisite)
\item \texttt{provides}, \texttt{eases}, \texttt{prefer}---the SPSE2 relations, preserved
\item \texttt{attacks}, \texttt{supports}---argumentation relations for contingent-plan reasoning
\item \texttt{subsumes}---HTN-style abstraction
\item \texttt{temporal/Allen-relation}---qualitative temporal constraints
\end{itemize}

Graph projection is composable. The projection predicate takes a filter
expression built from \texttt{kind/1}, \texttt{status/1},
\texttt{in\_mt/1}, \texttt{and/2}, \texttt{or/2}, \texttt{not/1}, and
returns the induced subgraph:

\begin{lstlisting}
% Project the dependency subgraph of medical_kb,
% excluding completed tasks:
?- project_graph(medical_kb,
                 and(kind(depends), not(status(completed))),
                 proj(Nodes, Edges)).
\end{lstlisting}

The SPSE2 legacy importer reads \texttt{holds/2} dumps and maps them to
the new ontology:

\begin{lstlisting}
?- import_spse2_file('~/spse2-dumps/pse.holds', legacy_pse).
\end{lstlisting}

The importer translates \texttt{holds(Ctx, goal('entry-fn'(P,I)))} to a
task, \texttt{holds(Ctx, 'has-NL'(\ldots, NL))} to the \texttt{has\_nl}
string, \texttt{holds(Ctx, depends(A,B))} to an \texttt{edge(A, depends,
B)}, \texttt{holds(Ctx, complete(\ldots))} to \texttt{status=completed},
and the other SPSE2 statuses, \texttt{costs}, \texttt{earns}, etc.\
analogously.

\subsection{\texttt{pack-pddl}: PDDL/HDDL Parser and Emitter}
\label{sec:pddl}

Approximately 560 lines, 20 tests. A full DCG-based parser and emitter for
PDDL 1.2, 2.1 \citep{fox-long-pddl21-2003}, 2.2 \citep{edelkamp-pddl22-2004},
3.0, 3.1, and HDDL, together with a
planner-capability table and a \texttt{required\_features/2} utility that
inspects a parsed domain for the PDDL features it uses.

Key predicates: \texttt{parse\_pddl/2} (parse a PDDL file into a structured
Prolog term); \texttt{emit\_pddl/2} (emit a structured term as PDDL text);
\texttt{required\_features/2} (inspect a parsed domain, return the set of
PDDL features it uses, such as \texttt{:typing},
\texttt{:durative-actions}, \texttt{:timed-initial-literals});
\texttt{eligible\_planners/2} (given a feature set, return planners capable
of handling it, from the capability table).

The round-trip property \emph{parse $\circ$ emit $\circ$ parse $\equiv$
parse} is verified by a dedicated test. The planner table covers the
common temporal and HTN planners and is data-driven---adding a new planner
is a single \texttt{planner\_capability/2} fact.

\textbf{Why this matters standalone.} No other SWI-Prolog pack provides a
complete PDDL parser across the language versions. The feature-detection
utility is particularly useful for the autopackager's workflow of deciding
which planner to hand a given domain to.

\subsection{\texttt{pack-spse4-scheduler}: Real-Datetime Scheduling and Reactive Triggers}
\label{sec:scheduler}

Approximately 450 lines, 15 tests. Extends \texttt{pack-allen}'s
integer-timeline scheduler to real calendar time, adds deadline propagation
and critical-path analysis, and wires \texttt{when/2}-driven reactive
triggers that fire when a task's dependencies become satisfied.

Real datetimes use SWI-Prolog's built-in \texttt{library(date)} (Unix-epoch
seconds as the internal representation) with \texttt{format\_time/3} for
display. No external date library is required; callers who prefer the
\texttt{julian} pack can still use its values as inputs since any integer
seconds-since-epoch is accepted.

Core capabilities:

\begin{itemize}
\item \texttt{schedule\_tasks/3}---schedule all tasks in a microtheory
  respecting asserted temporal relations, durations, earliest\_starts, and
  latest\_finishes.
\item \texttt{critical\_path/3}---compute the critical path from ready
  tasks to a goal task, returning the path along with its length in
  seconds.
\item \texttt{propagate\_deadlines/2}---given a \texttt{latest\_finish} on
  a downstream task and known durations of its upstream dependencies,
  compute the implicit \texttt{latest\_finish} of the upstream tasks
  (classical PERT backward pass).
\item \texttt{task\_earliest\_start/3}, \texttt{task\_latest\_finish/3},
  \texttt{task\_slack/3}---point-query interface for individual task timing
  and slack (latest-start minus earliest-start).
\item \texttt{on\_dependencies\_complete/3}---register a goal that should
  fire when a given task's dependencies all reach \texttt{status=completed}.
\item \texttt{trigger\_ready\_tasks/1}---scan a microtheory and broadcast
  \texttt{ready\_task\_event(Mt, TaskId)} for every task whose dependencies
  are now satisfied.
\end{itemize}

The reactive layer uses \texttt{library(broadcast)} rather than
\texttt{when/2} directly, because \texttt{when/2}'s coroutining semantics
tie the trigger to a specific query scope; broadcast gives the
server-authoritative fanout the v0.2 collaboration layer will need (see
\citet{lager-pengines-2014} for the Pengines transport this enables).

% ==========================================================================
\section{The Autopackager Demo: An End-to-End Walkthrough}
\label{sec:demo}
% ==========================================================================

The demo at \texttt{autopackager\_demo.pl} seeds SPSE4 with eleven tasks
drawn from the FRDCSA autopackager's actual priority queue---package
\texttt{eprover}, \texttt{peleus}, \texttt{vampire}, \texttt{acl2},
\texttt{freekbs2}, submit to mentors, prepare the FLP public release, and
several intermediate milestones. It then exercises every major SPSE4
capability on real data.

\subsection{Seeding the microtheory}

The demo opens by creating a microtheory and asserting eleven tasks with
their dependencies, durations, statuses, and declared earliest\_starts.
Three are already \texttt{completed} (\texttt{pkg\_eprover},
\texttt{pkg\_peleus}, \texttt{pkg\_superlu}), one is \texttt{in\_progress}
(\texttt{pkg\_vampire}), and the rest are \texttt{open}.

\begin{lstlisting}
Seeded autopackager microtheory with 11 tasks.
\end{lstlisting}

\subsection{Identifying ready tasks}

The demo then projects the dependency graph and identifies which tasks
have all prerequisites satisfied, producing the list of tasks that could
be started immediately.

\subsection{Critical path to the FLP public release}

\texttt{critical\_path/3} is called with \texttt{pkg\_flp\_public\_release}
as the goal, and returns the critical path along with its cumulative
duration. On the seeded data this path runs through \texttt{pkg\_vampire}
(the current in-progress task), which is the actual bottleneck in the
real queue.

\subsection{Schedule projection with real dates}

\texttt{schedule\_tasks/3} runs the full CLP(FD) schedule against the
asserted durations and earliest\_starts. For the tasks with declared start
dates (\texttt{pkg\_vampire} at April 22, 2026), the scheduler places them
on their declared dates; for tasks without declared starts, it schedules
them starting from ``today'' (the demo's run date). Dependencies are
honored: \texttt{submit\_autopkg\_to\_mentors} is placed after
\texttt{pkg\_vampire} completes.

Representative output (edited for brevity):

\begin{lstlisting}
pkg_vampire        : 2026-04-22 -> 2026-04-25  (3 days)
submit_mentors     : 2026-04-25 -> 2026-04-26  (1 day, waits on vampire)
pkg_acl2           : 2026-04-22 -> 2026-04-24
pkg_freekbs2       : 2026-04-22 -> 2026-04-23
pkg_flp_release    : 2026-05-04 -> 2026-05-04  (milestone)
\end{lstlisting}

\subsection{Readiness gate with blocker identification}

A gate check on \texttt{submit\_autopkg\_to\_mentors} reports:

\begin{lstlisting}
  pkg_eprover: completed
  pkg_peleus: completed
  pkg_vampire: in_progress
  NOT YET -- 1 blocker(s) outstanding:
    pkg_vampire (in_progress)
\end{lstlisting}

The gate not only returns a boolean but names the specific blockers and
their current statuses---this is the actual information a user needs to
act on the gate failure.

\subsection{Reactive trigger demonstration}

Finally the demo registers a reactive trigger: ``when
\texttt{submit\_autopkg\_to\_mentors} becomes ready, announce so.'' The
trigger registers silently. The demo then simulates \texttt{pkg\_vampire}
completing (by updating its status to \texttt{completed} and calling
\texttt{trigger\_ready\_tasks/1}), and the trigger fires automatically
with the announcement.

\begin{lstlisting}
Simulating vampire completion:
  *** AUTOMATIC: all prereqs complete, time to submit! ***
\end{lstlisting}

\subsection{Why this demo matters}

Every piece of the classical command-and-control
observe-orient-decide-act loop is represented in this single end-to-end
run: perception (task status queries), comprehension (dependency graph
projection), projection (the schedule), decision (gate checks,
critical-path identification), and action (reactive trigger fires on
state change). This is real C2 reasoning on real operational data---not a
toy. It is the first time SPSE4 has done something meaningful with the
author's actual planning queue, and it validates the architecture: the
five packs compose cleanly, the demo is $\approx$200 lines of application
code, and the user-visible behavior is correct on the first full
end-to-end run.

% ==========================================================================
\section{Running the System}
% ==========================================================================

\subsection{Installation}

Install SWI-Prolog $\geq$ 9.0.0. Then either (a) use the five packs as a
meta-project from the SPSE4 source tree, or (b) install each pack
individually via \texttt{pack\_install/1} from its GitHub URL (planned for
v0.2).

Meta-project use:

\begin{lstlisting}
cd ~/Downloads/spse4
swipl run_all_tests.pl          # expected: 79 tests green, no warnings
swipl autopackager_demo.pl      # expected: end-to-end demo output
\end{lstlisting}

\subsection{Loading individual packs from Prolog}

\begin{lstlisting}
?- use_module(library(allen)).
?- use_module(library(mt_store)).
?- use_module(library(spse4_core)).
?- use_module(library(pddl)).
?- use_module(library(spse4_scheduler)).
\end{lstlisting}

\subsection{Per-pack test suites}

Each pack ships a \texttt{t/*.plt} PlUnit file. Run with:

\begin{lstlisting}
swipl -g "[t/allen], run_tests" -t halt pack-allen/prolog/allen.pl
\end{lstlisting}

% ==========================================================================
\section{Coding Conventions}
% ==========================================================================

The entire codebase follows a small set of explicit rules, consistent
across all five packs:

\begin{itemize}
\item \emph{Covington-Bratko style.} Lowercase-underscore predicate names,
  \texttt{:- module/2} on every source file, PlDoc docstrings on every
  exported predicate, PlUnit tests colocated as \texttt{t/<module>.plt}.
\item \emph{Strings for user-visible text.} Labels, descriptions,
  natural-language fragments, free-form notes---all SWI strings
  (\texttt{""} quoted).
\item \emph{Atoms for identifiers and enums.} Task IDs, relation kinds,
  status values, microtheory names, predicate names---all Prolog atoms
  (\texttt{''} or unquoted).
\item \emph{No side effects outside designated I/O modules.} Predicates
  that mutate the store are in \texttt{mt\_store} and nowhere else.
\item \emph{Every exported predicate has both a ground-case test and a
  variable-binding test.} This catches the most common class of bug (a
  predicate that works for queries but fails when asked to generate
  bindings, or vice versa).
\end{itemize}

These rules came out of the SPSE4 v0.1 session and are recorded here
because they are the main reason the codebase held together across twelve
bug-fix iterations without structural refactoring.

% ==========================================================================
\section{Roadmap to v0.2}
\label{sec:roadmap}
% ==========================================================================

v0.1 delivers the foundational layer. The next release delivers the
collaboration layer and two additional reasoning layers.

\subsection{\texttt{pack-spse4-server}: Pengines HTTP/WebSocket Server}

A SWI-Prolog Pengines server with basic authentication, per-microtheory
access control enforcement, and \texttt{library(broadcast)} fanout so
clients see diffs in near-real-time. Approximately 300 lines. Endpoints
follow the standard Pengines API \citep{lager-pengines-2014}
(\texttt{pengine/create},
\texttt{pengine/send}, \texttt{pengine/pull\_response},
\texttt{pengine/destroy}) extended with an SPSE4-specific
projection-filter endpoint that returns only the nodes and edges visible
under a given filter expression.

\subsection{\texttt{spse4-web}: Cytoscape.js Client}

A browser client using Cytoscape.js for graph rendering and
\texttt{pengines.js} for the Prolog connection. The projection UI exposes
checkboxes for \texttt{relation\_kind}, \texttt{in\_mt}, and
\texttt{status}; nodes are colored by status and by critical-path
membership. Approximately 500 lines of JS/HTML, single page, no framework.

\subsection{\texttt{spse4-mode.el}: Emacs Mode}

An Emacs mode using \texttt{pengine.el} (already written) to provide the
same projection operations from inside Emacs, so SPSE4 becomes directly
editable from \texttt{org-mode} buffers without leaving the editor.

\subsection{Additional reasoning layers}

Three additional packs are planned for v0.2 or v0.3, each independently
motivated:

\begin{itemize}
\item \emph{Native Prolog behavior-tree executor.} Reactive execution of
  task decompositions, with tick-based evaluation and standard BT control
  flow (sequence, selector, parallel, decorator). No external runtime
  dependency.
\item \emph{Event Calculus layer (RTEC-inspired).} For temporal reasoning
  beyond Allen intervals, particularly for recurrence and for reasoning
  about what happened when; \texttt{pack-allen} answers ``what qualitative
  relation holds'', the Event Calculus layer will answer ``what fluent
  holds at time $t$''.
\item \emph{s(CASP) plan justification.} For every scheduled plan, produce
  a proof tree showing why each scheduling decision was forced; lets the
  user audit the scheduler's reasoning step-by-step.
\end{itemize}

\subsection{Exports}

\begin{itemize}
\item \emph{CycL text emission.} Already sketched (see \S3.5); v0.2 will
  finalize the emitter and add a round-trip test using a stable
  \texttt{.cycl} corpus.
\item \emph{RDF/Turtle export.} For federation with Koo5's work and for
  compatibility with external knowledge-graph tooling.
\item \emph{HTN SHOP-style meta-interpreter.} Task decomposition using the
  SHOP HTN style, with SPSE4's \texttt{subsumes} edges as the
  method-precondition source.
\item \emph{planutils adapter.} A thin adapter for the \texttt{planutils}
  polyglot-planner backend, so \texttt{emit\_pddl} $\to$ \texttt{planutils}
  $\to$ parsed plan is a one-liner.
\end{itemize}

% ==========================================================================
\section{Discussion}
% ==========================================================================

Three properties of the v0.1 delivery are worth highlighting explicitly.

\textbf{Iteration discipline.} The session produced twelve tarballs
(v0.1.0 $\to$ v0.1.12). Each tarball fixed exactly one bug, caught with
one targeted diagnostic, with no structural overshoot. The bugs fell into
categories that are well-documented SWI-Prolog gotchas---\texttt{assertz}
copying attribute variables, \texttt{forall/2} undoing CLP(FD) bindings,
choice-point warnings on tail-recursive helpers---and each fix was local.
The root cause, in every case, was a subtle semantics mismatch between the
programmer's model and the actual Prolog semantics; none of the bugs were
design flaws. This is consistent with the ``tight feedback loop on real
data'' methodology the FRDCSA project has consistently used.

\textbf{Scope discipline.} v0.1 is aggressively scoped. Things
deliberately deferred include: the collaboration server, the browser
client, the Event Calculus layer, the s(CASP) justification layer, live
Cyc integration, and the MySQL backend for \texttt{mt\_store}. Every
deferral has a stated rationale. This is what allowed v0.1 to ship with
79 green tests and an end-to-end working demo in a single session.

\textbf{Dogfooding from day one.} The autopackager demo uses the author's
actual packaging queue, not a toy. This matters for two reasons. First,
it forces the user-visible ergonomics to be correct---there is no way to
fake a realistic demo. Second, it means SPSE4 is a working tool for the
author on the day of the v0.1 release, not a tool waiting for content to
be useful. The FRDCSA project has consistently prioritized this kind of
self-hosting (``the Free Life Planner must first plan its own
development''), and SPSE4 v0.1 meets the bar.

% ==========================================================================
\section*{Acknowledgments}
\addcontentsline{toc}{section}{Acknowledgments}
% ==========================================================================

SPSE4 builds on intellectual lineage from \citet{allen-1983} on
interval algebra, \citet{guha-lenat-cyc-1990} on microtheories and
context logic, the PDDL community
\citep{fox-long-pddl21-2003,edelkamp-pddl22-2004}, and the SWI-Prolog
core team \citep{wielemaker-swipl-2012}. The 2011 SPSE2 technical
report provided the ontology. The v0.1 codebase was developed in a
live pair-programming session between the author and Claude
(Anthropic, Claude Opus 4.6 and 4.7) on April 22, 2026, with Claude
producing the majority of the Prolog source and the author providing
domain knowledge, architectural judgment, and test-failure
diagnostics. The iteration pattern of ``paste the error, fix one
thing, re-run'' proved remarkably effective across the twelve
tarballs.

% ==========================================================================
\section*{License}
\addcontentsline{toc}{section}{License}
% ==========================================================================

All five packs ship under the MIT License. The project is hosted on
GitHub under the \texttt{aindilis} organization at
\url{https://github.com/aindilis/} (repos \texttt{pack-allen},
\texttt{pack-mt-store}, \texttt{pack-pddl}, \texttt{pack-spse4-core},
\texttt{pack-spse4-scheduler}, and the \texttt{spse4} umbrella).

\bibliographystyle{plainnat}
\bibliography{spse4}
\end{document}
